Take one admissible design, let \(X\) be constant, let \(U\) be uniform on \(\{1,2,3\}\), and let \(S_2\) be an independent fair bit, with the same distribution in both sources. In source \(E\), set \(P(A=1\mid G=E)=1/2\). In source \(O\), set
\[
 P(A=1\mid U=u,G=O)=e_u,\qquad (e_1,e_2,e_3)=(1/5,1/2,4/5).
\]
Conditional on \((S_2=s,U=u)\), let \(S_1\) and \(S_3\) be independent Bernoulli variables, unaffected by treatment and source, with common success probability \(p_{s,u}\), where
\[
 (p_{0,1},p_{0,2},p_{0,3})=(1/5,1/2,4/5),\qquad
 (p_{1,1},p_{1,2},p_{1,3})=(3/10,4/5,2/5).
\]
Use the same \(S_3\) in both potential short-term systems and set
\[
 Y(0)=S_3,\qquad Y(1)=S_3+\delta_{S_2}(U),
\]
where
\[
 (\delta_0(1),\delta_0(2),\delta_0(3))=(4,-16/5,1),\qquad
 (\delta_1(1),\delta_1(2),\delta_1(3))=(-16/5,-8/25,1).
\]
For each \(s\in\{0,1\}\), direct multiplication gives
\[
 \sum_{u=1}^3 e_u\delta_s(u)=0,
 \qquad
 \sum_{u=1}^3 e_up_{s,u}\delta_s(u)=0.                 \tag{1}
\]
The two equations in (1), corresponding to \(S_1=0\) and \(S_1=1\), show that
\[
 h_d(s_3,s_2,a)=s_3
\]
satisfies the observed outcome-bridge equation for both arms. Thus the required square-integrable observed bridge exists.

For completeness, the matrix whose rows are \(1\), \(p_0\), and \(p_1\) has determinant \(-27/100\). For either treatment arm, conditioning on that arm merely multiplies its three columns by a positive diagonal matrix. Hence
\[
 \mathbb E\{v(U)\mid S_1,S_2,A=a,G=O\}=0
 \quad\Longrightarrow\quad v(U)=0.
\]
The identical argument with \(S_3\) proves late-proxy completeness. For each \((s,a)\), the \(2\)-by-\(2\) conditional-expectation matrix from functions of \(S_3\) to functions of \(S_1\) has determinant proportional to the strictly positive conditional variance of \(p_{s,U}\); it is therefore nonsingular. Thus \(T_d\) is bijective. In finite dimension \(T_d^\star\) is also bijective, so the observed selection-bridge equation has a unique square-integrable solution.

All source and arm probabilities are strictly positive. Conditional independence of \(S_1\) and \(S_3\) given \((S_2,U)\) supplies the sequential restriction. Consistency, observational latent exchangeability, experimental randomization, external validity, and source support follow directly from the construction. Finite support gives finite moments. The experimental score is nondegenerate because \(h_d=S_3\), and the observational score is nondegenerate because the treated residual contains the nonzero \(\delta_{S_2}(U)\). Consequently the constructed law belongs to \(\mathfrak P_{\mathcal D}\).

The distribution of \(S_3\) in the randomized source is treatment invariant, so
\[
 \theta_d(P)=
 \mathbb E\{S_3\mid A=1,G=E\}
 -\mathbb E\{S_3\mid A=0,G=E\}=0.
\]
In contrast,
\[
 \tau
 =\mathbb E\{\delta_{S_2}(U)\}
 =\frac16\left(\frac95-\frac{63}{25}\right)
 =-\frac3{25}.
\]
This proves the claimed separation and shows that common identification is an additional condition.